\chapter{项目二：正运动学、逆运动学与雅可比}

\section{项目二要解决的完整问题}

项目二的输入是目标末端位姿
\[
\T_{\mathrm{target}}=
\begin{bmatrix}
\R_{\mathrm{target}}&\p_{\mathrm{target}}\\
0&1
\end{bmatrix},
\]
输出是若干组可能的关节角。它不是一个孤立函数，而是一条验证链：
\[
\q\xrightarrow{\mathrm{FK}}\T
\xrightarrow{\mathrm{IK}}\q'
\xrightarrow{\mathrm{FK}}\T'.
\]
只有当 $\T'\approx\T$ 时，逆解才可信。

\section{SR3 的 MDH 参数}

\begin{table}[H]
\centering
\caption{当前代码采用的 SR3 MDH 参数，长度单位为 mm}
\begin{tabular}{ccccc}
\toprule
$i$ & $a_{i-1}$ & $\alpha_{i-1}$ & $d_i$ & $\theta_i$\\
\midrule
1 & 0   & $0$             & 344   & $q_1+\pi$\\
2 & 0   & $\pi/2$         & 0     & $q_2+\pi/2$\\
3 & 290 & $\pi$           & 0     & $q_3+\pi/2$\\
4 & 0   & $\pi/2$         & 290   & $q_4+\pi$\\
5 & 0   & $-\pi/2$        & 0     & $q_5$\\
6 & 136 & $\pi/2$         & 103.5 & $q_6$\\
\bottomrule
\end{tabular}
\end{table}

\section{正运动学如何一步一步计算}

\subsection{第一步：把关节角变成 MDH 表}

\code{build\_mdh\_table.m} 对每个关节执行
\[
\theta_i=q_i+\mathrm{offset}_i,
\]
然后把 $[a_{i-1},\alpha_{i-1},d_i,\theta_i]$ 放入第 $i$ 行。

\subsection{第二步：每一行变成一个局部变换}

调用
\[
\boldsymbol{A}_i=
\mathrm{mdh\_transform}(a_{i-1},\alpha_{i-1},d_i,\theta_i).
\]
$\boldsymbol{A}_i$ 只描述“从前一坐标系走到当前坐标系”。

\subsection{第三步：把六步串起来}

\[
{}^0\T_6=
{}^0\boldsymbol{A}_1
{}^1\boldsymbol{A}_2
{}^2\boldsymbol{A}_3
{}^3\boldsymbol{A}_4
{}^4\boldsymbol{A}_5
{}^5\boldsymbol{A}_6.
\]
代码从 \code{T = eye(4)} 开始，每次右乘一个 $\boldsymbol{A}_i$。
同时保存中间结果 \code{T\_all(:,:,i+1)}，因为碰撞检测和雅可比不仅需要
末端，还需要每个关节原点的位置。

\section{逆运动学为什么比正运动学难}

正运动学只需按确定顺序代入并相乘，通常只有一个结果。逆运动学困难在于：

\begin{itemize}
  \item 同一末端位姿可能对应多组关节角；
  \item 有些目标超出工作空间，根本无解；
  \item 接近奇异位形时，某些关节运动效果重合；
  \item 三角函数是周期函数，$q$ 与 $q+2k\pi$ 表示相同转角；
  \item 机械臂还受关节限位约束。
\end{itemize}

\section{课程要求的 Pieper 分离法}

\subsection{核心思想}

对后三根关节轴交于一点的六轴球腕机械臂，可以把问题拆开：

\begin{enumerate}
  \item 前三轴负责把腕中心送到正确位置；
  \item 后三轴负责把末端转到正确朝向。
\end{enumerate}

因此先解位置，再解姿态。这称为位置与姿态解耦。

\subsection{腕中心}

在最简单的球腕模型中，末端原点沿末端 $z$ 轴距离腕中心 $d_6$，因此
\[
\p_w=\p_6-d_6\R_6(:,3).
\]
$\R_6(:,3)$ 是末端 $z$ 轴在基坐标系中的方向。

\begin{dangerbox}[当前 SR3 参数并非最简单模板]
当前模型第六行同时含 $a_5=136$ 和 $d_6=103.5$。这意味着末端固定偏置不只
沿 $z_6$ 方向，简单公式 $\p_w=\p-d_6\R(:,3)$ 不足以完整剥离偏置。
这也是当前仓库最终转向通用数值 IK 的重要原因。不能把模板球腕公式不加检查地
宣称为当前代码的精确实现。
\end{dangerbox}

\subsection{前三关节的几何解}

对于理想化平面二连杆，设腕中心为
$[x_w,y_w,z_w]^T$。第一轴先决定腕中心位于机器人哪一侧：
\[
\theta_1=\operatorname{atan2}(y_w,x_w).
\]
另一肩部构型常可写为 $\theta_1+\pi$。

把三维问题投影到关节 2、3 所在平面：
\[
r=x_w\cos\theta_1+y_w\sin\theta_1,\qquad
s=z_w-d_1.
\]
若两段有效长度为 $L_1,L_2$，根据余弦定理：
\[
D=\frac{r^2+s^2-L_1^2-L_2^2}{2L_1L_2}.
\]
当 $|D|>1$ 时三角形无法组成，目标不可达。当 $|D|\le1$ 时：
\[
\theta_3=
\operatorname{atan2}\left(\pm\sqrt{1-D^2},D\right).
\]
正负号产生肘上、肘下两组解。再由
\[
\theta_2=
\operatorname{atan2}(s,r)
-\operatorname{atan2}
\left(L_2\sin\theta_3,L_1+L_2\cos\theta_3\right)
\]
求得肩关节俯仰角。

\subsection{后三关节的姿态解}

前三关节确定后，先算 $\R_{03}$，目标姿态为 $\R_{06}$，于是
\[
\R_{36}=\R_{03}^{-1}\R_{06}=\R_{03}^{T}\R_{06}.
\]
再把 $\R_{36}$ 与腕部三个旋转矩阵的乘积逐元素比较，就能通过
\code{atan2} 提取 $\theta_4,\theta_5,\theta_6$。

一般情况下 $\theta_5$ 的正负分支对应腕部翻转，从而形成两组腕部解。
于是最多有
\[
2\text{（肩）}\times2\text{（肘）}\times2\text{（腕）}=8
\]
组解。

\subsection{腕部奇异}

当 $\sin\theta_5\approx0$ 时，第 4 轴和第 6 轴的效果耦合。此时不是没有解，
而是 $\theta_4$ 与 $\theta_6$ 各自不唯一，通常只有两者之和或差能够确定。
程序必须人为选一个约定，例如令 $\theta_4=0$，再计算 $\theta_6$。

\section{当前代码真正使用的 DLS 数值逆运动学}

\subsection{把逆运动学改写成“不断纠错”}

假设当前猜测为 $\q_k$。用正运动学得到当前位姿 $\T(\q_k)$，与目标比较得到
六维误差：
\[
\e=
\begin{bmatrix}
\e_p\\
\e_o
\end{bmatrix}.
\]
前三维 $\e_p$ 是位置误差：
\[
\e_p=\p_{\mathrm{target}}-\p_{\mathrm{current}}.
\]
后三维 $\e_o$ 是姿态误差。代码先算
\[
\R_{\mathrm{err}}=
\R_{\mathrm{target}}\R_{\mathrm{current}}^T,
\]
再把误差旋转矩阵转换为“旋转轴乘旋转角”的三维向量。

\subsection{小运动近似与雅可比}

在关节角变化很小时：
\[
\Delta\boldsymbol{x}\approx\J(\q)\Delta\q.
\]
希望 $\Delta\boldsymbol{x}$ 抵消当前误差，即
\[
\J\Delta\q\approx\e.
\]
如果直接使用 $\Delta\q=\J^{-1}\e$，会遇到两个问题：$\J$ 可能不可逆，
接近奇异位形时数值会爆炸。

\subsection{阻尼最小二乘}

DLS 使用
\[
\Delta\q=
\J^T\left(\J\J^T+\lambda^2\boldsymbol{I}\right)^{-1}\e.
\]
其中 $\lambda>0$ 是阻尼系数。$\lambda^2I$ 像给病态矩阵加一层缓冲，
使其更容易求解。

代码没有写 \code{inv(A)}，而是：
\begin{lstlisting}
A = J * J' + (opts.lambda^2) * eye(6);
dq = J' * (A \ e);
\end{lstlisting}
因为 \code{A\textbackslash e} 会选择更稳定的线性方程求解方法。

更新规则是
\[
\q_{k+1}=\q_k+\Delta\q^T.
\]
迭代直到 $\norm{\e}$ 小于阈值，或达到最大迭代次数。

\subsection{阻尼系数的影响}

\begin{itemize}
  \item $\lambda$ 太小：远离奇异点时收敛快，但奇异附近可能出现巨大步长；
  \item $\lambda$ 太大：更稳定，但每次只走很小一步，收敛慢且可能残留误差；
  \item 当前单次 IK 默认常用 $0.01$，多初值封装中使用 $0.05$。
\end{itemize}

\section{几何雅可比如何推导}

\subsection{雅可比的尺寸}

六自由度机械臂的几何雅可比为
\[
\J=
\begin{bmatrix}
\J_v\\
\J_\omega
\end{bmatrix}\in\mathbb{R}^{6\times6}.
\]
它把关节速度映射为末端速度：
\[
\begin{bmatrix}\boldsymbol{v}\\\boldsymbol{\omega}\end{bmatrix}
=\J(\q)\dot{\q}.
\]

\subsection{转动关节的一列}

第 $i$ 个转动关节的轴方向为 $\boldsymbol{z}_i$，轴上一点为 $\p_i$，
末端为 $\p_n$。该关节单位角速度引起：
\[
\J_{v,i}=\boldsymbol{z}_i\times(\p_n-\p_i),\qquad
\J_{\omega,i}=\boldsymbol{z}_i.
\]

\subsection{MDH 中为什么不能直接拿累计矩阵的 z 轴}

本项目单步顺序为
\[
\R_x(\alpha)\T_x(a)\R_z(\theta)\T_z(d).
\]
真正的关节转轴是执行前两个固定变换之后、执行 $\R_z(\theta)$ 之前的
$z$ 轴。所以代码先计算
\[
\T_{\mathrm{pre}}=
\T\,\R_x(\alpha_i)\T_x(a_i),
\]
再从 \code{T\_pre(1:3,3)} 取轴方向，从
\code{T\_pre(1:3,4)} 取轴上一点。

\section{姿态误差的轴角表示}

任意旋转矩阵对应一个旋转轴 $\boldsymbol{u}$ 和角度 $\theta$。角度由
\[
\cos\theta=\frac{\operatorname{tr}(\R)-1}{2}
\]
得到。旋转轴乘角度为
\[
\boldsymbol{\omega}=
\frac{\theta}{2\sin\theta}
\begin{bmatrix}
R_{32}-R_{23}\\
R_{13}-R_{31}\\
R_{21}-R_{12}
\end{bmatrix}.
\]
当 $\theta$ 很小时，分母也很小，因此代码直接返回零向量，避免数值除零。

\section{多初值为什么能找到多组解}

DLS 是局部方法：从哪个初始点出发，通常会落入附近的某一个解。
\code{build\_seed\_set} 先准备若干人为设计的肩、肘、腕构型，再在关节限位中
生成 30 组固定随机种子。每个种子独立运行 DLS，最终可能到达不同解。

候选解需经过四道筛选：

\begin{enumerate}
  \item 数值迭代必须声明收敛；
  \item 正运动学回代的位置和姿态误差必须足够小；
  \item 必须满足关节限位；
  \item 与已有解的周期角距离大于阈值，否则视为重复解。
\end{enumerate}

\section{角度归一化与解的选择}

角度具有周期性：
\[
q\equiv q+2k\pi.
\]
代码用
\[
\operatorname{wrap}(x)=\operatorname{mod}(x+\pi,2\pi)-\pi
\]
把角度差折算到 $[-\pi,\pi)$。

选解函数先过滤关节限位，再选择离参考构型最近者：
\[
k^*=\arg\min_k
\norm{\operatorname{wrap}(\q^{(k)}-\q_{\mathrm{ref}})}_2.
\]
这样做的工程意义是减少关节突然大幅跳动，保持动作连续。

\section{当前实现的优点与局限}

\subsection{优点}

\begin{itemize}
  \item 不依赖严格球腕结构，可直接适配当前完整 MDH 模型；
  \item 每个解都经过 FK 反校验，结果可信；
  \item DLS 在奇异附近比普通伪逆稳定；
  \item 多初值可以得到多组候选解。
\end{itemize}

\subsection{局限}

\begin{itemize}
  \item 不能数学保证找到全部解，也不能严格保证最多恰好八组；
  \item 计算量明显高于闭式解析解；
  \item 收敛依赖初值、阻尼和迭代上限；
  \item 位置以 mm、姿态以 rad 混在同一误差向量中，未设置尺度权重；
  \item 接近旋转角 $\pi$ 时，简单轴角公式也可能数值敏感。
\end{itemize}

\begin{warningbox}[报告答辩建议]
不要说“代码已经完整实现 Pieper 闭式解析法”。更准确的表述是：
“报告学习并推导了 Pieper 解耦思想；由于当前 SR3 MDH 末端偏置与简化球腕模板
不完全一致，工程代码采用多初值 DLS，并用 FK 严格反校验候选解。”
\end{warningbox}
